% Per-game replacements of the fixed texts of lang-es.tex (\renewcommand).
\renewcommand{\tCoverLine}{El juego de MSX turbo R en el V9990}
\renewcommand{\tMemoryText}{\textbf{\ramNeeded{} de RAM}; todos los MSX turbo R la tienen.}
\renewcommand{\tVideoText}{Un \textbf{cartucho V9990} (vídeo Yamaha V9990). Es obligatorio: toda la imagen la dibuja el V9990. Salvo los mensajes del encendido, el juego no usa la pantalla propia del ordenador.}
\renewcommand{\tSoundText}{Conecta los altavoces al \textbf{MoonSound o Wozblaster} para oír juntos la música y los efectos. El cartucho mezcla la música FM con los disparos, explosiones y demás sonidos del arcade en PCM. Con un cartucho que sólo tenga OPL3, los efectos PSG salen por el audio del MSX y la música por el cartucho. Sin cartucho de sonido, usa la salida del MSX para los efectos.}
\renewcommand{\tSetGame}{Inserta el \textbf{cartucho de \gameTitle{}} en otra ranura: un slot primario o un subslot de un expansor de slots. El orden no importa.}
\renewcommand{\tSetOn}{Enciende el MSX. Un momento la pantalla del propio MSX muestra \texttt{V9990 found !}; después la imagen aparece en el monitor del V9990. Si dice \texttt{THIS GAME RUNS ONLY ON MSX TURBO R} el ordenador no es un turbo R, y si dice \texttt{V9990 NOT found!!!} falta el V9990 o no está bien insertado; en los dos casos el juego se detiene ahí.}
\renewcommand{\tControlsText}{Teclado (teclas del cursor) o un joystick de dos botones en el puerto 1.}
\renewcommand{\tPowerOnCaption}{Izquierda: la portada. Derecha: la pantalla de título con las siete mejores puntuaciones.}
\renewcommand{\tCreditsTitle}{10. Créditos y derechos}
\newcommand{\tCreditsExtra}{\item La imagen de portada es la ilustración de las versiones domésticas del juego original. La historia de la sección~5 resume fuentes públicas (Wikipedia y otras).}
